\documentclass[a4paper]{article}

%------------------------------------------------
% الحزم
%------------------------------------------------

\usepackage{enumitem}
\usepackage[margin=2.2cm]{geometry}
\usepackage{amsmath,amssymb}
\usepackage{graphicx}
\usepackage{array}
\usepackage{fancyhdr}
\usepackage[table]{xcolor}
\usepackage{tikz}
\usetikzlibrary{positioning,arrows.meta,shapes.geometric,calc}
\usepackage[skins]{tcolorbox}
\usepackage{polyglossia}

%------------------------------------------------
% اللغات والخطوط
%------------------------------------------------

\setdefaultlanguage[
    calendar=gregorian,
    numerals=maghrib
]{arabic}

\setotherlanguage{english}

% إذا لم تكن الخطوط مثبتة في الجهاز، تُؤخذ من المجلد fonts/
\IfFontExistsTF{Amiri}{%
    \newfontfamily\arabicfont[Script=Arabic, Scale=1.15]{Amiri}
}{%
    \newfontfamily\arabicfont[
        Script=Arabic, Scale=1.15,
        Path=fonts/, Extension=.ttf,
        UprightFont=*-Regular, BoldFont=*-Bold,
        ItalicFont=*-Italic, BoldItalicFont=*-BoldItalic
    ]{Amiri}
}

\IfFontExistsTF{AlBattar}{%
    \newfontfamily\arabicfontsf[Script=Arabic, Scale=1.2]{AlBattar}
    % خط العناوين (يعمل حتى بعد استعمال \textenglish)
    \newfontfamily\khatt[Script=Arabic, Scale=1.2]{AlBattar}
}{%
    \newfontfamily\arabicfontsf[
        Script=Arabic, Scale=1.2, Path=fonts/,
        BoldFont=ae_AlBattar.ttf
    ]{ae_AlBattar.ttf}
    % خط العناوين (يعمل حتى بعد استعمال \textenglish)
    \newfontfamily\khatt[
        Script=Arabic, Scale=1.2, Path=fonts/,
        BoldFont=ae_AlBattar.ttf
    ]{ae_AlBattar.ttf}
}

%------------------------------------------------
% الألوان: لون رئيسي هادئ + لون تنبيه فقط
%------------------------------------------------

\definecolor{asasi}{HTML}{3D6A8E}   % اللون الرئيسي: للرسوم والمخططات
\definecolor{tanbih}{HTML}{C0603A}  % لون التنبيه: للتحذيرات والحالات السلبية
\definecolor{ramadi}{HTML}{5F6B73}  % رمادي للنصوص الثانوية والأسهم
\definecolor{dakin}{HTML}{263238}   % لون العناوين

% أسماء قديمة تُحوَّل إلى اللونين فقط
\colorlet{azraq}{asasi}
\colorlet{fayruzi}{asasi}
\colorlet{banafsaji}{asasi}
\colorlet{akhdar}{asasi}
\colorlet{ahmar}{tanbih}
\colorlet{burtuqali}{asasi}

%------------------------------------------------
% الصناديق
%------------------------------------------------

\tcbset{
    asl/.style={
        enhanced,
        arc=2mm,
        boxrule=0.6pt,
        colback=white,
        colframe=black!18,
        colbacktitle=black!5,
        coltitle=dakin,
        titlerule=0pt,
        left=8pt, right=8pt, top=5pt, bottom=5pt,
        fonttitle=\khatt\bfseries,
        before skip=8pt, after skip=8pt
    }
}

% أهداف الدرس
\newtcolorbox{ahdaf}{asl, title={أهداف الدرس}}

% تعريف (خط رئيسي رفيع على الجانب)
\newtcolorbox{taarif}[1][تعريف]{asl, title={#1},
    borderline east={2pt}{0pt}{asasi}}

% مثال
\newtcolorbox{mithal}[1][مثال]{asl, title={#1}}

% تنبيه (اللون الوحيد البارز في النص)
\newtcolorbox{intabih}[1][انتبه!]{asl, colback=tanbih!5, colframe=tanbih,
    colbacktitle=tanbih, coltitle=white, title={#1}}

% ركن أستاذ اللغة العربية
\newtcolorbox{arabiyya}[1][في قسم اللغة العربية]{asl, title={#1}}

% نص أمر يكتب للأداة
\newtcolorbox{amr}[1][ramadi]{enhanced, arc=1mm, boxrule=0pt,
    colback=black!4, borderline east={2.5pt}{0pt}{#1},
    left=8pt, right=10pt, top=4pt, bottom=4pt,
    before skip=6pt, after skip=6pt}

% الخلاصة
\newtcolorbox{khulasa}{asl, title={خلاصة الدرس}}

% عنوان القسم
\newcounter{qism}
\newcommand{\qism}[1]{%
    \stepcounter{qism}%
    \par\bigskip
    {\large\khatt\bfseries{\color{asasi}\arabic{qism}.}\ {\color{dakin}#1}}\par
    \vspace{-2pt}{\color{black!20}\rule{\linewidth}{0.6pt}}\par\nopagebreak[4]
    \smallskip}

% عنوان فرعي
\newcommand{\far}[1]{\par\smallskip
    {\color{dakin}\khatt\bfseries {\color{asasi}$\blacktriangleleft$}\ #1}\par\nopagebreak[4]}

% كلمة مهمة (خط عريض فقط)
\newcommand{\kalima}[2][]{\textbf{#2}}

%------------------------------------------------
% أنماط الرسوم
%------------------------------------------------

\newcommand{\nodecolor}{}

\tikzset{
    % كتابة النص العربي من اليمين إلى اليسار داخل العقد
    % (لون النص يُعطى بالمفتاح lawn=... وليس text=...)
    every node/.append style={
        execute at begin node={\nodecolor\everypar{\beginR}\beginR},
        execute at end node=\endR},
    lawn/.code={\def\nodecolor{\color{#1}}},
    khatwa/.style={draw=#1, fill=#1!12, very thick, rounded corners=6pt,
        minimum height=1.2cm, text width=2.3cm, align=center,
        font=\khatt\bfseries},
    khatwa/.default=azraq,
    sahm/.style={-{Stealth[length=3mm]}, very thick, color=ramadi},
    raqm/.style={circle, fill=#1, lawn=white, inner sep=2pt,
        font=\small\bfseries},
}

%------------------------------------------------
% إعداد الصفحة
%------------------------------------------------

\renewcommand{\footrulewidth}{1pt}
\renewcommand{\headrulewidth}{1pt}

\rhead{\khatt المدرسة العليا للأساتذة ببوزريعة}
\lhead{\khatt ملحقة بومرداس}

\cfoot{\thepage}
\lfoot{\khatt الأستاذ: د. بن عبد الرزاق يوسف}

\parindent 0pt
\setlength{\headsep}{15pt}
\setlength{\parskip}{4pt}

\pagestyle{fancy}

\setlist[enumerate,1]{label=\arabic*), itemsep=2pt, topsep=3pt}
\setlist[itemize]{itemsep=2pt, topsep=3pt}

% الجداول
\arrayrulecolor{black!22}
\renewcommand{\arraystretch}{1.35}

% رأس جدول ملون
\newcommand{\ras}[1]{\cellcolor{black!7}\color{dakin}\bfseries #1}

% عنوان الدرس
\newcommand{\unwan}[3]{%
    \begin{center}
        {\color{asasi}\rule{\linewidth}{1.2pt}}\par\medskip
        {\khatt\color{ramadi} #1}\par\smallskip
        {\Large\khatt\bfseries\color{dakin} #2}\par\smallskip
        {\khatt\color{ramadi} #3}\par\medskip
        {\color{asasi}\rule{\linewidth}{1.2pt}}
    \end{center}}

%================================================
% بداية الوثيقة
%================================================

\begin{document}

%------------------------------------------------
% العنوان
%------------------------------------------------

\unwan{درس الفصل السابع}
{تصميم مساعد ذكاء اصطناعي وسير عمل}
{تعليمات ــ معرفة ــ اختبار ــ تحسين ــ سير عمل}

\begin{ahdaf}
في نهاية هذا الدرس، تستطيع أن:
\begin{itemize}
    \item تشرح ما المساعد المخصص، وما مكوناته.
    \item تكتب تعليمات واضحة لمساعد تعليمي.
    \item تختبر المساعد وتحسّن تعليماته.
    \item تقسّم مهمة كبيرة إلى سير عمل من عدة مراحل.
\end{itemize}
\end{ahdaf}

%================================================
\qism{لماذا نبني مساعدًا مخصصًا؟}
%================================================

يستعمل الأستاذ الذكاء الاصطناعي كل أسبوع: يلخّص الدروس، ويشرح المفاهيم، ويُعدّ الأسئلة.
لكنه يكتب التعليمات نفسها في كل محادثة جديدة، ويرفع الملفات نفسها في كل مرة.

الحل: نبني \kalima{مساعدًا مخصصًا} مرة واحدة، ثم نستعمله كلما احتجنا إليه.

\begin{taarif}
\textbf{المساعد المخصص:} نسخة من أداة الذكاء الاصطناعي،
نعطيها اسمًا وتعليمات ثابتة وملفات مرجعية، لتقوم بمهمة محددة.
\end{taarif}

نبنيه في \textenglish{ChatGPT} باستعمال \textenglish{GPTs}،
أو في \textenglish{Claude} باستعمال \textenglish{Projects}.

%================================================
\qism{مكونات المساعد}
%================================================

\begin{center}
\begin{tikzpicture}[
    juz/.style={draw=#1, fill=#1!12, very thick, rounded corners=6pt,
        text width=3.3cm, align=center, minimum height=1.6cm}]
    \node[circle, draw=azraq, fill=azraq!20, ultra thick,
        text width=2.3cm, align=center, font=\khatt\bfseries] (m) {المساعد};
    \node[juz=akhdar] (a) at (4.3,1.3)
        {{\khatt\bfseries الاسم}\\\small مساعد مادة النحو};
    \node[juz=burtuqali] (b) at (4.3,-1.3)
        {{\khatt\bfseries التعليمات}\\\small \textenglish{Instructions}};
    \node[juz=banafsaji] (c) at (-4.3,1.3)
        {{\khatt\bfseries المعرفة}\\\small \textenglish{Knowledge}: ملفات المادة};
    \node[juz=fayruzi] (d) at (-4.3,-1.3)
        {{\khatt\bfseries الأدوات}\\\small البحث، تحليل البيانات...};
    \foreach \x in {a,b,c,d} \draw[very thick, ramadi] (m) -- (\x);
\end{tikzpicture}
\end{center}

\begin{center}
\rowcolors{2}{black!3}{white}
\begin{tabular}{>{\centering\arraybackslash}m{2.6cm}
                >{\raggedleft\arraybackslash}m{11.3cm}}
\hline
\ras{المكوّن} & \ras{دوره} \\
\hline
\kalima[akhdar]{الاسم}         & يعرّف المساعد ومهمته من أول نظرة. \\
\kalima[burtuqali]{التعليمات}  & تحدد دور المساعد وطريقة عمله وحدوده. \\
\kalima[banafsaji]{المعرفة}    & ملفات مرجعية يعتمد عليها: دروس، ملخصات، محاضرات. \\
\kalima[fayruzi]{الأدوات}      & قدرات إضافية. لا نفعّلها إلا إذا احتاجها المساعد فعلًا. \\
\hline
\end{tabular}
\end{center}

%================================================
\qism{كتابة التعليمات}
%================================================

التعليمات هي «عقل» المساعد.
نكتبها باستعمال عناصر الأمر الجيد التي درسناها في الفصل الثالث:

\begin{amr}[ramadi]
\kalima{الدور:} أنت مساعد تعليمي متخصص في مادة النحو العربي.\\
\kalima[akhdar]{المهمة:} ساعد الطلبة على فهم الدروس بطريقة بسيطة ومنظمة.\\
\kalima[fayruzi]{طريقة العمل:} اشرح خطوة بخطوة، واستعمل أمثلة بسيطة.\\
\kalima[banafsaji]{المراجع:} اعتمد على الملفات المرفقة عندما يتعلق السؤال بالمادة.\\
\kalima[ahmar]{القيود:} لا تخترع معلومات. إذا لم تجد المعلومة، صرّح بذلك.\\
\kalima[burtuqali]{شكل الإجابة:} عناوين واضحة، ونقاط قصيرة، وأمثلة.\\
\kalima[ramadi]{السؤال الغامض:} اطلب توضيحًا قبل أن تجيب.
\end{amr}

%================================================
\qism{دورة بناء المساعد}
%================================================

بناء المساعد لا ينتهي بإنشائه.
نختبره، ونكتشف أخطاءه، ونحسّنه، ثم نختبره من جديد:

\begin{center}
\begin{tikzpicture}[
    marhala/.style={draw=#1, fill=#1!14, very thick, rounded corners=6pt,
        text width=2.2cm, align=center, minimum height=1cm,
        font=\khatt\bfseries}]
    \node[marhala=azraq]     (a) at (90:3)  {حدّد المهمة};
    \node[marhala=fayruzi]   (b) at (30:3.4)   {اكتب التعليمات};
    \node[marhala=akhdar]    (c) at (-30:3.4)  {أنشئ وأضف الملفات};
    \node[marhala=burtuqali] (d) at (-90:3) {اختبر};
    \node[marhala=ahmar]     (e) at (-150:3.4) {اكتشف الأخطاء};
    \node[marhala=banafsaji] (f) at (150:3.4)  {حسّن التعليمات};
    \draw[sahm] (a) to[bend left=20] (b);
    \draw[sahm] (b) to[bend left=20] (c);
    \draw[sahm] (c) to[bend left=20] (d);
    \draw[sahm] (d) to[bend left=20] (e);
    \draw[sahm] (e) to[bend left=20] (f);
    \draw[sahm, dashed] (f) to[bend left=10] node[above, sloped, font=\small, lawn=ramadi]
        {أعد الاختبار} (d);
\end{tikzpicture}
\end{center}

%================================================
\qism{اختبار المساعد}
%================================================

لا نفترض أن المساعد يعمل جيدًا. نختبره بأربعة أنواع من الأسئلة:

\begin{center}
\rowcolors{2}{black!3}{white}
\begin{tabular}{>{\centering\arraybackslash}m{3cm}
                >{\raggedleft\arraybackslash}m{5.6cm}
                >{\raggedleft\arraybackslash}m{5cm}}
\hline
\ras{الاختبار} & \ras{مثال} & \ras{ماذا نلاحظ؟} \\
\hline
سؤال بسيط             & ما المبتدأ؟ اشرحه لطالب مبتدئ.  & هل اللغة بسيطة؟ \\
سؤال يعتمد على الملف   & اعتمد على الملف، واذكر أهم ثلاثة مفاهيم. & هل الإجابة مطابقة للملف؟ \\
سؤال ناقص             & اشرح لي هذا الموضوع بالتفصيل.     & هل يطلب توضيحًا؟ \\
سؤال خارج المهمة       & كيف أحضّر البيتزا؟                & هل يحافظ على دوره؟ \\
\hline
\end{tabular}
\end{center}

\begin{mithal}[تحسين بعد الاختبار]
لاحظنا أن المساعد يعطي إجابات طويلة جدًا. نضيف إلى التعليمات:
\begin{amr}[akhdar]
استعمل جملًا قصيرة ونقاطًا واضحة.
لا تقدم شرحًا طويلًا إلا إذا طلب المستخدم ذلك.
\end{amr}
ثم نعيد الاختبار نفسه، ونقارن الإجابة قبل التعديل وبعده.
\end{mithal}

%================================================
\qism{سير العمل \textenglish{(Workflow)}}
%================================================

\begin{taarif}
\textbf{سير العمل:} تقسيم مهمة كبيرة إلى مراحل صغيرة متتابعة.
مخرج كل مرحلة يصبح مدخلًا للمرحلة التالية.
\end{taarif}

\far{مثال: من درس طويل إلى أسئلة تدريبية}

\begin{center}
\begin{tikzpicture}[node distance=4mm and 5mm]
    \node[khatwa=ramadi, text width=1.9cm] (a) {النص\\الأصلي};
    \node[khatwa=azraq, text width=1.9cm, left=of a] (b) {استخراج\\الأفكار};
    \node[khatwa=akhdar, text width=1.9cm, left=of b] (c) {إنشاء\\الملخص};
    \node[khatwa=burtuqali, text width=1.9cm, left=of c] (d) {إنشاء\\الأسئلة};
    \node[khatwa=banafsaji, text width=1.9cm, left=of d] (e) {التحقق\\النهائي};
    \draw[sahm] (a)--(b); \draw[sahm] (b)--(c);
    \draw[sahm] (c)--(d); \draw[sahm] (d)--(e);
    % نقاط المراجعة
    \foreach \x in {b,c,d} {
        \coordinate (o) at ($(\x.south)+(0,-0.55)$);
        \draw[ahmar, very thick, fill=white] (o) circle (0.22);
        \draw[ahmar, very thick] ($(o)+(-0.11,0)$) -- ($(o)+(-0.03,-0.1)$)
            -- ($(o)+(0.12,0.1)$);
    }
    \node[font=\small, lawn=ahmar] at ($(c.south)+(0,-1.25)$)
        {مراجعة بشرية بعد كل مرحلة};
\end{tikzpicture}
\end{center}

\begin{center}
\begin{tabular}{>{\centering\arraybackslash}m{1.6cm}
                >{\centering\arraybackslash}m{3.6cm}
                >{\centering\arraybackslash}m{4cm}
                >{\centering\arraybackslash}m{3.6cm}}
\hline
\ras{المرحلة} & \ras{المدخل} & \ras{المعالجة} & \ras{المخرج} \\
\hline
1 & نص المحاضرة & استخراج الأفكار الأساسية & قائمة أفكار \\ \hline
2 & قائمة الأفكار & كتابة ملخص منظم & ملخص الدرس \\ \hline
3 & الملخص & إنشاء خمسة أسئلة مع الإجابات & أسئلة تدريبية \\ \hline
\end{tabular}
\end{center}

\begin{intabih}
لا ننتقل إلى المرحلة التالية قبل مراجعة نتيجة المرحلة الحالية.
الخطأ الصغير في البداية يكبر في المراحل التالية.
\end{intabih}

\far{أمر واحد أم سير عمل؟}

\begin{center}
\begin{tabular}{>{\raggedleft\arraybackslash}m{5cm}
                >{\centering\arraybackslash}m{4cm}
                >{\centering\arraybackslash}m{4cm}}
\hline
\ras{المعيار} & \ras{أمر واحد} & \ras{سير عمل} \\
\hline
سهولة المراجعة         & \kalima[ahmar]{صعبة}   & \kalima[akhdar]{سهلة} \\ \hline
اكتشاف الأخطاء         & \kalima[ahmar]{صعب}    & \kalima[akhdar]{سهل} \\ \hline
التحكم في كل مرحلة     & \kalima[ahmar]{ضعيف}   & \kalima[akhdar]{قوي} \\ \hline
إعادة الاستعمال        & \kalima[burtuqali]{محدودة} & \kalima[akhdar]{ممكنة} \\ \hline
السرعة                 & \kalima[akhdar]{أسرع}  & \kalima[burtuqali]{أبطأ قليلًا} \\ \hline
\end{tabular}
\end{center}

\begin{arabiyya}[أفكار لمساعدين في قسم اللغة العربية]
\begin{itemize}
    \item \textbf{مساعد الإعراب:} يشرح إعراب الجمل خطوة بخطوة، اعتمادًا على كتاب النحو المقرر.
    \item \textbf{مساعد البلاغة:} يستخرج الصور البيانية من نص، ويشرحها.
    \item \textbf{مساعد التعبير:} يقترح ملاحظات على إنشاء التلميذ دون أن يكتبه بدلًا منه.
\end{itemize}
\end{arabiyya}

\begin{khulasa}
\begin{itemize}
    \item المساعد المخصص = \textbf{اسم} + \textbf{تعليمات} + \textbf{معرفة} + \textbf{أدوات}.
    \item نبني، ثم \textbf{نختبر}، ثم \textbf{نحسّن}، ثم نختبر من جديد.
    \item \textbf{سير العمل} يقسم المهمة الكبيرة إلى مراحل، مع مراجعة بعد كل مرحلة.
    \item المساعد الجيد يحتاج إلى أستاذ يراجع نتائجه باستمرار.
\end{itemize}
\end{khulasa}

\end{document}
